\documentclass[a4j,10.5pt,dvipdfmx]{jsreport}

% 数式
\usepackage{amsmath,amsfonts}
\usepackage{bm}

% 画像
\usepackage[dvipdfmx,top=35truemm,bottom=30truemm,left=30truemm,right=30truemm]{geometry}
\usepackage[dvipdfmx]{graphicx}
% biblatex
\usepackage[style=numeric, sorting=none, maxbibnames=1000]{biblatex}
\addbibresource{graduation_thesis.bib}

% itembox
\usepackage{ascmac}
% breakitemboxのためのpackage
\usepackage{itembkbx}
% \usepackage{tcolorbox} 画像が見えなくなってしまう原因
% \usepackage{framed} 枠

% フォント
\usepackage[ipaex]{pxchfon}
\listfiles
\usepackage{float}
\usepackage{multicol}
\usepackage{multirow}
\usepackage{paralist}
\usepackage{caption}
% \captionsetup{format=plain, justification=centering}
\usepackage{setspace}
\usepackage{algorithm}
\usepackage{algpseudocode}
\usepackage[dvipdfmx,hidelinks]{hyperref}
\usepackage{pxjahyper}
\usepackage{fancyhdr}
\usepackage{ltablex}
\usepackage{capt-of}
\usepackage{tabularx}
\usepackage[export]{adjustbox}
\usepackage{booktabs}
\usepackage{enumitem} % itemizeの行間 これは下になければならない
\usepackage{threeparttable} % 表の脚注
\usepackage{titlesec}
\usepackage{makecell}
\usepackage{forest}
\usepackage{fancybox}
% titlesec と \paragraph の相性問題を回避（run-in 見出しにする）
\titleformat{\paragraph}[runin]
  {\normalfont\normalsize\bfseries}
  {}{0pt}{} % ←ここには文字や \noindent 等を入れない
\titlespacing*{\paragraph}
  {0pt}{1.0ex}{0.8em}
\titleformat{\chapter}[hang]{\LARGE\bfseries}{\prechaptername\thechapter\postchaptername}{1em}{}

\setlength{\headheight}{16.5pt}
\captionsetup{format=hang}
\captionsetup[figure]{font=small}
\captionsetup[table]{font=small}
\begin{document}

% ▼▼▼ 仕掛け：main じゃないとき（単体のとき）だけヘッダー設定 ▼▼▼
\ifdefined\ismain\else
   \pagestyle{fancy}
   \lhead{\rightmark}
   \renewcommand{\chaptermark}[1]{\markboth{第\ \thechapter\ 章\ #1}{}}
   \setlength{\headheight}{16.5pt}
   \titleformat{\chapter}[hang]{\LARGE\bfseries}{\prechaptername\thechapter\postchaptername}{1em}{}
\fi
% ▲▲▲▲▲▲▲▲▲▲▲▲▲▲▲▲▲▲▲▲▲▲▲▲▲▲▲▲▲▲▲▲▲▲▲▲

\begin{abstract}
   現在，AIは創作領域に新たな可能性を開きつつある一方，倫理的課題を抱え，社会的受容には依然として隔たりがある．とりわけ，既存作品の安易な模倣に基づく粗悪な派生物を拡散する危険性を持ち，これは日本が培ってきた創作資産の価値と流通を毀損する恐れがある．日本のコンテンツ産業を守り持続的な発展を維持するためにも，その価値が適切に還元され得るAI活用の枠組みを整備する必要がある．そこで本研究では，最終成果物そのものではなく，制作工程で用いられるレイアウトのような中間的生成物に着目し，創作を支援するAI活用の可能性を検討する．とりわけ漫画制作には，コマ割りやセリフ配置などの演出上の意思決定を担う内部資料として，ネームが存在する．
   漫画は平面をコマへと分割し，その上に物語を描く表現形式である．コマやキャラクター，セリフの配置は演出上の意思決定を反映し，一般にネームと呼ばれる下書きにおいて定められ，作品の核となる役割を担う．一方，要素の配置を扱うレイアウト生成の研究は主に機能的媒体を対象として発展してきたため，ネームのように物語文脈と不可分に決まる配置を対象とした研究は十分ではない．さらにネームは内部資料として公開されにくく，学習・評価に耐える規模で脚本情報や演出情報が整備されたデータも不足していた．
   そこで本研究では，ネームを下書き画像そのものではなく，要素の幾何学的配置と演出情報が統合された，物語から作品へ至るための構造化中間表現として定義する．この定義に基づき，Manga109Dataset に含まれる19,555ページの漫画画像から脚本情報および演出情報を抽出し，Manga109ScriptDataset と Manga109NameDatasetとして整備した．この過程では，漫画読解の要となる読み順を明示するために，103,850のコマおよび127,887のセリフに順序を人手で付与し，これを視覚的手がかりとしたビジュアルプロンプティング手法により情報抽出精度を高めた．さらに，これらのデータセットを用いて，脚本を入力としネームを生成する漫画ネーム生成モデルを構築した．提案モデルは既存のレイアウト生成手法との比較により，複数の評価指標で性能の向上を示した．その一方で，演出情報の定量的な評価指標の不足が課題となった．以上により，本研究は，漫画制作の中核であるネームを対象とした大規模データの整備と，脚本に基づくネーム生成の学習基盤を提示する．
\end{abstract}


\newpage % 必要であれば改ページ
% 英語の要旨タイトルに変更
\renewcommand{\abstractname}{\Huge Thesis Abstract}
% 英語の要旨
\begin{quotation}
   AI creates new possibilities for art, but ethical issues remain. Low-quality copies made by AI can hurt the value of Japan’s creative assets. To support the industry, we explore AI that helps with the process rather than just making final outputs. We focus on the "name" (storyboard), a key document in manga that decides panel layouts and dialogue placement.
   Manga tells a story by splitting a page into panels. This layout reflects the author's intent. However, **previous layout generation research** focused on simple media and ignored layouts tied to stories. Also, because storyboards are private, there is not enough data for machine learning.
   In this study, we define a storyboard as a structured representation combining layout and direction. We extracted data from 19,555 pages to create the \textit{Manga109ScriptDataset} and \textit{Manga109NameDataset}. To ensure accuracy, we manually annotated the reading order of **103,850 panels and 127,887 dialogue balloons**. We used this order as a visual cue in a visual prompting method to improve data extraction. Using these datasets, we built a model that generates storyboards directly from scripts. Our model performs better than existing methods on several metrics, though quantitatively evaluating "directing" quality remains a challenge. Overall, this work provides large-scale resources and a foundation for manga storyboard generation.
\end{quotation}


% ▼▼▼ 仕掛け：main じゃないとき（単体のとき）だけ参考文献を表示 ▼▼▼
\ifdefined\ismain\else
   \printbibliography[title=参考文献]
\fi
% ▲▲▲▲▲▲▲▲▲▲▲▲▲▲▲▲▲▲▲▲▲▲▲▲▲▲▲▲▲▲▲▲▲▲▲▲

\end{document}